\icode{E-MK-SI-L2-260928-002}
公路上行驶的汽车，司机从发现前方异常情况到紧急刹车，汽车仍将前进一段距离才能停下来。要保证安全，这段距离内不能有车辆和行人，因此把它称为安全距离。由于人和车的灵敏程度不同，导致反应时间不同，一般需要\(0.8\sim\)\qty{1.0}{s}（在做出反应前，汽车仍然匀速行驶）。由于车况不同，刹车后汽车获得的制动力不同，一般是车重的\(0.5\sim0.6\)倍。重力加速度\(g =\)\qty{10}{m/s}在限速\qty{108}{km/h}的高速公路上，交警要对所有司机提示一个安全车距，请你通过计算告诉交警该安全车距至少应定为\paren[D]
\begin{choices}
	\item \qty{99}{m}
	\item \qty{105}{m}
	\item \qty{114}{m}
	\item \qty{120}{m}
\end{choices}
\begin{answer}
	D\qquad 提示：取最大速度、最小制动力和最长反应时间计算距离
\end{answer}